\documentclass[11pt]{article}

\usepackage[a4paper,margin=2.5cm]{geometry}
\usepackage{amsmath,amssymb,amsthm}
\usepackage{booktabs}
\usepackage{array}
\usepackage[colorlinks=true,linkcolor=blue]{hyperref}

\newtheorem{lemma}{Lemma}
\renewcommand{\thelemma}{A.\arabic{lemma}}
\theoremstyle{remark}
\newtheorem{remark}{Remark}
\renewcommand{\theremark}{A.\arabic{remark}}

\newcommand{\Deff}{D_{\mathrm{eff}}}
\newcommand{\dcol}{d_{\mathrm{col}}}
\newcommand{\Twall}{T_{\mathrm{wall}}}
\newcommand{\lameff}{\lambda_{\mathrm{eff}}}
\newcommand{\DTad}{\Delta T_{\mathrm{ad}}}

\title{Mechanistic Fixed-Bed CO$_2$ Adsorption Model\\[2pt]
\large Part A --- Derivation from Conservation Laws}
\author{}
\date{}

\begin{document}
\maketitle

\section*{Part A --- Derivation from conservation laws}

\subsection*{A.1 Gas-phase CO$_2$ mass balance}

Take a control volume $[z,z+\Delta z]$ of cross-section $A$ (bed basis). Moles of
gaseous CO$_2$ inside: $\varepsilon\,c\,A\,\Delta z$.

\paragraph{Fluxes through the faces} (per bed area, $\mathrm{mol\,m^{-2}\,s^{-1}}$):
\begin{equation*}
N(z,t) \;=\; \underbrace{u\,c}_{\text{convective}}\;
\underbrace{-\;\varepsilon D_L \frac{\partial c}{\partial z}}_{\text{dispersive (Fickian)}}.
\end{equation*}
The convective term is $u c = \varepsilon v_i c$: interstitial molar flux $v_i c$
scaled by the open-area fraction $\varepsilon$. The dispersive flux acts on the same
open area, hence the $\varepsilon$.

\paragraph{Sink to the solid.} Sorbent mass in the CV is $(1-\varepsilon)\rho_p A\Delta z$,
so uptake removes $(1-\varepsilon)\rho_p \frac{\partial q}{\partial t} A\Delta z$
$\mathrm{mol\,s^{-1}}$ from the gas.

Balance (accumulation $=$ in $-$ out $-$ sink), divide by $A\Delta z$, let $\Delta z\to0$:
\begin{equation*}
\varepsilon\frac{\partial c}{\partial t} + \frac{\partial N}{\partial z}
= -(1-\varepsilon)\rho_p\frac{\partial q}{\partial t},
\end{equation*}
and with A2 ($u_z=0$):
\begin{equation}
\boxed{\;\varepsilon\frac{\partial c}{\partial t} + u\frac{\partial c}{\partial z}
\;=\; \varepsilon D_L\frac{\partial^2 c}{\partial z^2}
\;-\;(1-\varepsilon)\rho_p\frac{\partial q}{\partial t}\;}
\tag{A.1}\label{eq:A1}
\end{equation}

\begin{center}
\begin{tabular}{@{}l p{8cm} l@{}}
\toprule
Term & Physical meaning & Units \\
\midrule
$\varepsilon\,c_t$ & accumulation of gaseous CO$_2$ per \textbf{bed} volume
  ($\varepsilon$ converts gas-basis $c$ to bed basis) & $\mathrm{mol\,m^{-3}\,s^{-1}}$ \\
$u\,c_z$ & net convective outflow per bed volume, superficial $u$ & $\mathrm{mol\,m^{-3}\,s^{-1}}$ \\
$\varepsilon D_L c_{zz}$ & net dispersive inflow (axial mixing) & $\mathrm{mol\,m^{-3}\,s^{-1}}$ \\
$(1-\varepsilon)\rho_p q_t = \alpha_b q_t$ & transfer to the adsorbed phase & $\mathrm{mol\,m^{-3}\,s^{-1}}$ \\
\bottomrule
\end{tabular}
\end{center}

\begin{remark}[variable velocity --- when 10--15\,\% feeds bite]\label{rem:A1}
Only CO$_2$ leaves the gas; a total molar balance with $c_{\mathrm{tot}}=P/RT$
constant (A3, isothermal) gives
\begin{equation*}
c_{\mathrm{tot}}\frac{\partial u}{\partial z} = -(1-\varepsilon)\rho_p\frac{\partial q}{\partial t}
\quad\Longrightarrow\quad
u(z,t)=u_{\mathrm{in}}-\frac{\alpha_b}{c_{\mathrm{tot}}}\int_0^z q_t\,dz'.
\end{equation*}
Across a saturating front the velocity deficit is $\Delta u/u \approx y_f$:
${\sim}1\,\%$ at $1\,\%$ feed (negligible) but ${\sim}10$--$15\,\%$ for runs 4--8
(\texttt{new runs/}). The dilute model \eqref{eq:A1} then mislocates $t_{st}$ by
$O(y_f)$; if that exceeds the $\pm10\,\%$ Gate-B budget, solve \eqref{eq:A1}
together with the $u$-equation above (still MOL-compatible: $u$ is obtained by
quadrature of $q_t$ at each RHS evaluation).
\end{remark}

\subsection*{A.2 Solid-phase balance: LDF kinetics}

Exact pellet-scale mass transfer is a diffusion problem in the particle; the
\textbf{linear driving force} closure replaces it by first-order relaxation toward
equilibrium:
\begin{equation}
\boxed{\;\frac{\partial q}{\partial t} = k\,\bigl(q^*(c,T)-q\bigr)\;}
\tag{A.2}\label{eq:A2}
\end{equation}

\begin{center}
\begin{tabular}{@{}l l l@{}}
\toprule
Term & Physical meaning & Units \\
\midrule
$q_t$ & uptake rate per kg sorbent & $\mathrm{mol\,kg^{-1}\,s^{-1}}$ \\
$k$ & inverse relaxation time of the pellet & $\mathrm{s^{-1}}$ \\
$q^*-q$ & displacement from equilibrium (driving force) & $\mathrm{mol\,kg^{-1}}$ \\
\bottomrule
\end{tabular}
\end{center}

\paragraph{Mechanistic content of $k$} (resistances in series, Glueckauf; Ruthven
ch.~6): for spherical pellets of radius $r_p$, film coefficient $k_f$, macropore
diffusivity $D_p$, porosity $\varepsilon_p$, crystal/micro scale $r_c$, $D_c$, and
local slope $K'=\rho_p\,\partial q^*/\partial c$,
\begin{equation*}
\frac{1}{k} \;=\; \frac{r_p K'}{3k_f} \;+\; \frac{r_p^2 K'}{15\,\varepsilon_p D_p}
\;+\; \frac{r_c^2}{15\,D_c},
\end{equation*}
each term a time constant (s). LDF is exact for a linear isotherm with parabolic
intraparticle profile; otherwise it is the controlled approximation this project
fits at Gate~B. For amine sorbents (PEI@SiO$_2$) the last, reaction/amine-phase
term often dominates and $k$ is small --- the measured, strongly smeared fronts
(\texttt{breakthrough\_out/run 5}: $t_b\approx14$~s vs $t_{50}\approx230$~s) are
consistent with a kinetics-limited bed (\S C.4).

\subsection*{A.3 Equilibrium closure}

\paragraph{Toth (primary; amine-functionalised sorbent).} Written in the
concentration basis,
\begin{equation}
\boxed{\;q^*(c,T)=\frac{n_s(T)\,b(T)\,c}{\bigl[1+(b(T)\,c)^{t_T}\bigr]^{1/t_T}}\;}
\tag{A.3}\label{eq:A3}
\end{equation}
\begin{equation*}
n_s(T)=n_{s0}\exp\!\Bigl[\chi\Bigl(1-\tfrac{T}{T_0}\Bigr)\Bigr],\qquad
b(T)=b_0\exp\!\Bigl[\tfrac{\Delta H_0}{R T_0}\Bigl(\tfrac{T_0}{T}-1\Bigr)\Bigr],\qquad
t_T(T)=t_0+\alpha_T\Bigl(1-\tfrac{T_0}{T}\Bigr),
\end{equation*}
with $0<t_T\le 1$ the heterogeneity exponent ($t_T=1$ recovers Langmuir).
\textbf{Units trap:} literature Toth parameters (e.g.\ the Stampi-Bombelli 2024
set, still \texttt{??} in \texttt{CLAUDE.md}) are usually \emph{pressure-basis},
$b_P$ in $\mathrm{kPa^{-1}}$ with $q^*(p_{\mathrm{CO_2}})$; convert with
$p_{\mathrm{CO_2}}=cRT$, i.e.\ $b(T)=b_P(T)\,R\,T$. The conversion inserts an
extra $T$-dependence --- do it before differentiating, not after.

Equivalent driving-force form used in \texttt{derivation.md} \S1.2 ($C-C^*$ with
$C^*$ the isotherm inverse): identical model when
$k_a a_p (1-\varepsilon)(C-C^*) = \alpha_b k (q^*-q)$ under the local
linearisation $k = k_a a_p (C-C^*)/(\rho_p (q^*-q))$; \eqref{eq:A2} is preferred
because $q^*(c,T)$ stays single-valued while the inverse $C^*(q,T)$ need not be
evaluated.

\paragraph{Dual-Site Langmuir (alternative; zeolite 13X):}
\begin{equation}
q^*(c,T)=\frac{q_1 b(T)c}{1+b(T)c}+\frac{q_2 d(T)c}{1+d(T)c},\qquad
b=b_0e^{-\Delta U_b/RT},\quad d=d_0e^{-\Delta U_d/RT},\quad
\Delta U_b,\Delta U_d<0.
\tag{A.4}\label{eq:A4}
\end{equation}

\begin{lemma}[shape of the closures]\label{lem:A1}
For fixed $T$ and $c>0$, both \eqref{eq:A3} and \eqref{eq:A4} satisfy
\begin{enumerate}
\item $q^*(0,T)=0$ and Henry limit $q^*\sim n_s b\,c$ (Toth) resp.\
      $(q_1b+q_2d)c$ (DSL);
\item strict monotonicity --- Toth:
      $\dfrac{\partial q^*}{\partial c}
       = \dfrac{n_s b}{\bigl[1+(bc)^{t_T}\bigr]^{(1+t_T)/t_T}}>0$;
\item strict concavity --- Toth:
      $\dfrac{\partial^2 q^*}{\partial c^2}
       = -\,n_s b^2 (1+t_T)\,(bc)^{t_T-1}
         \bigl[1+(bc)^{t_T}\bigr]^{-(1+2t_T)/t_T}<0$;
      DSL: sum of Langmuir terms with $q_i b^2\cdot(-2)/(1+bc)^3<0$;
\item saturation $q^*\to n_s$ (Toth) resp.\ $q_1+q_2$ (DSL) as $c\to\infty$.
\end{enumerate}
\end{lemma}

\begin{proof}
Differentiate \eqref{eq:A3} with $s=(bc)^{t_T}$: $q^*=n_sbc(1+s)^{-1/t_T}$,
$\partial_c q^* = n_s b(1+s)^{-1/t_T}-n_s b s(1+s)^{-1/t_T-1}
= n_s b (1+s)^{-(1+t_T)/t_T}$; differentiate again using
$\partial_c s = t_T s/c$.
\end{proof}

Concavity $=$ \textbf{favorable} isotherm; it is what makes adsorption fronts
self-sharpening (Part~D). Note for $t_T<1$ the curvature blows up like
$c^{\,t_T-1}$ as $c\to0^+$: $q^*$ is $C^1$ but not $C^2$ at $c=0$ --- harmless
for well-posedness (the slope stays bounded by $n_sb$) but it steepens the
leading foot of the front.

\paragraph{A.4.3 Thermodynamic consistency of $(-\Delta H)$.}
The isosteric heat implied by the closure is
$(-\Delta H_{\mathrm{iso}}) = R T^2\,\partial_T \ln p \big|_{q^*}$. For
\eqref{eq:A3} with $\chi=\alpha_T=0$ this returns exactly
$-\Delta H_{\mathrm{iso}}=-\Delta H_0$, constant; with $\chi\neq0$ it becomes
loading-dependent. Assumption A7: use that same function in \eqref{eq:A5}.
Using an isotherm fitted with one $\Delta H_0$ and an energy balance with an
unrelated $(-\Delta H)$ violates the Clausius--Clapeyron relation and produces
spurious heat.

\subsection*{A.4 Pseudo-homogeneous energy balance}

Same control volume; energy carried by gas convection, conducted axially through
the composite bed, released by adsorption, exchanged with the wall. Accumulation
uses both phases (single $T$, A6):
\begin{equation}
\boxed{\;C_h\frac{\partial T}{\partial t}
+ u\rho_g c_{p,g}\frac{\partial T}{\partial z}
\;=\; \lameff\frac{\partial^2 T}{\partial z^2}
\;+\;(1-\varepsilon)\rho_p(-\Delta H)\frac{\partial q}{\partial t}
\;-\;\frac{4h_w}{\dcol}\,(T-\Twall)\;}
\tag{A.5}\label{eq:A5}
\end{equation}
\begin{equation*}
C_h=\varepsilon\rho_g c_{p,g}+(1-\varepsilon)\rho_p c_{p,s}
\quad[\mathrm{J\,m^{-3}\,K^{-1}}].
\end{equation*}

\begin{center}
\begin{tabular}{@{}l p{8.5cm} l@{}}
\toprule
Term & Physical meaning & Units \\
\midrule
$C_h T_t$ & sensible-heat accumulation of gas $+$ solid per bed volume & $\mathrm{W\,m^{-3}}$ \\
$u\rho_g c_{p,g} T_z$ & convected enthalpy gradient (superficial flux $u\rho_g c_{p,g}$,
  $\mathrm{J\,m^{-2}\,s^{-1}\,K^{-1}}$) & $\mathrm{W\,m^{-3}}$ \\
$\lameff T_{zz}$ & effective axial conduction/thermal dispersion & $\mathrm{W\,m^{-3}}$ \\
$\alpha_b(-\Delta H) q_t$ & adsorption heat source; $>0$ during uptake since
  $(-\Delta H)>0$, $q_t>0$ & $\mathrm{W\,m^{-3}}$ \\
$\tfrac{4h_w}{\dcol}(T-\Twall)$ & wall loss; $4/\dcol$ $=$ perimeter/area of the
  cylinder & $\mathrm{W\,m^{-3}}$ \\
\bottomrule
\end{tabular}
\end{center}

Neglected knowingly: adsorbed-phase heat capacity
$(1-\varepsilon)\rho_p q\,c_{p,a}$ ($\le$ a few \% of $C_h$ at
$q\lesssim1~\mathrm{mol\,kg^{-1}}$), pressure work, kinetic energy.
\textbf{Adiabatic column:} $h_w=0$.

\paragraph{Isothermal reduction --- quantitative justification (prompt item 3).}
Two independent smallness arguments:
\begin{enumerate}
\item \emph{Heat-release vs wall-loss timescales.} Heat is generated over the
front-passage time ${\sim}t_{st}$ and removed on $\tau_w = C_h \dcol/(4h_w)$.
The quasi-steady excursion is $\Delta T \approx \DTad\cdot \tau_w/t_{st}$ where
$\DTad=\alpha_b(-\Delta H)q_f/C_h$ is the adiabatic rise. Bench rig estimate
($\dcol=8.5$~mm, $C_h\sim6.6\times10^{5}~\mathrm{J\,m^{-3}\,K^{-1}}$,
$h_w\sim30~\mathrm{W\,m^{-2}\,K^{-1}}$, $t_{st}\sim600$~s):
$\tau_w\approx47$~s, $\DTad\approx40$~K, so $\Delta T\approx3$--$4$~K --- mild,
and the isothermal model is a defensible base case \textbf{for this rig}. For a
pilot column ($\dcol\gtrsim5$~cm) $\tau_w$ grows linearly in $\dcol$ and the
full \eqref{eq:A5} is mandatory.
\item \emph{Isotherm sensitivity.} The perturbation enters through $b(T)$:
$\delta q^*/q^* \sim (\Delta H_0/RT^2)\,\Delta T
\approx (70\,000/8.314/298^2)\times 3.5 \approx 0.33$\ldots{} i.e.\ even a
3--4~K excursion shifts local capacity by ${\sim}30\,\%$ near the front.
Conclusion: keep \eqref{eq:A5} in the model; drop it only after checking
\emph{both} numbers for the case at hand. (This is why the model here is
non-isothermal by default and the isothermal system is treated as a limit in
Part~D.)
\end{enumerate}

\subsection*{A.5 Initial and boundary conditions}

Clean bed:
\begin{equation}
c(z,0)=0,\qquad q(z,0)=0,\qquad T(z,0)=T_0.
\tag{A.6}\label{eq:A6}
\end{equation}

\paragraph{Inlet ($z=0$), feed step at $t>0$.} Two admissible choices:
\begin{itemize}
\item \emph{Dirichlet (hyperbolic-limit shortcut):} $c(0,t)=c_f$, $T(0,t)=T_f$.
\item \emph{Danckwerts flux continuity (correct with dispersion):} upstream of
the packing there is no dispersion, so the total flux $uc_f$ must equal the
total flux just inside:
\end{itemize}
\begin{equation}
u c_f = u\,c(0^+,t) - \varepsilon D_L\,c_z(0^+,t), \qquad
u\rho_g c_{p,g} T_f = u\rho_g c_{p,g} T(0^+,t) - \lameff T_z(0^+,t).
\tag{A.7}\label{eq:A7}
\end{equation}

\paragraph{Outlet ($z=L$):} $c_z(L,t)=0$, $T_z(L,t)=0$ (no dispersive flux
through the exit plane).

Danckwerts is a Robin condition; it makes the inventory identity of Part~B exact
and forces $c(0^+,t)<c_f$ while the front is inside the bed (the dispersive term
is positive). Dirichlet overfeeds the column by $-\varepsilon D_L c_z(0,t)>0$;
the defect is $O(Pe^{-1})$ and vanishes in the hyperbolic limit. Both are
implemented in Part~E; the finite-volume form makes Danckwerts trivial.

\subsection*{A.6 Full model (summary)}

\begin{equation*}
\varepsilon c_t + u c_z = \varepsilon D_L c_{zz} - \alpha_b q_t, \qquad
q_t = k\,(q^*(c,T)-q),
\end{equation*}
\begin{equation*}
C_h T_t + u\rho_g c_{p,g} T_z = \lameff T_{zz} + \alpha_b(-\Delta H)q_t
- \tfrac{4h_w}{\dcol}(T-\Twall),
\end{equation*}
with closure \eqref{eq:A3} or \eqref{eq:A4}, IC \eqref{eq:A6}, BC \eqref{eq:A7}
$+$ zero-gradient outlet. Structurally this is a
\textbf{parabolic--hyperbolic relaxation system}: two transport PDEs coupled
pointwise to one stiff local ODE.

\end{document}
